\documentclass[10pt,a4paper]{amsart}
\usepackage[margin=19mm]{geometry}
\usepackage{amsmath,amssymb,mathtools}
\usepackage[scr=rsfs,cal=euler]{mathalfa}
\usepackage{xfrac}
\usepackage[unicode,hidelinks,bookmarks=false]{hyperref}
\newcommand{\ie}{i.e.\ }
\newcommand{\cf}{cf.\ }
\newcommand{\resp}{respectively\ }
\newcommand{\Subsection}{\subsection}
\providecommand{\nobreakdash}{\nobreak\hbox{-}}
\setlength{\emergencystretch}{2em}
\multlinegap0pt
\usepackage{amssymb}
\newtheorem{theorem}{Theorem}
\newcommand{\CC}{\mathbb{C}}
\newcommand{\FF}{\mathbb{F}}
\newcommand{\ZZ}{\mathbb{Z}}
\newcommand{\supp}{\operatorname{supp}}
\newcommand{\tr}{\mathsf{T}}
\title{Tensor constructions for Euler magic matrices and proper examples of orders 9, 27, 81 and 243}
\author{Sanjit Singh Mehat}
\date{Source: arXiv:2609.05530v1 (CC BY 4.0)}
\begin{document}
\maketitle

\noindent\emph{Source and licence.} Tensor constructions for Euler magic matrices and proper examples of orders 9, 27, 81 and 243. Version arXiv:2609.05530v1 (CC BY 4.0), distributed by arXiv under the Creative Commons Attribution 4.0 International licence, http://creativecommons.org/licenses/by/4.0/, as recorded in the arXiv licence metadata for this exact version and shown on the abstract page https://arxiv.org/abs/2609.05530v1. Full licence notice: \texttt{licenses/arxiv-2609.05530-CC-BY-4.0.txt}.\par\smallskip

\noindent\textit{Editorial scope.} Counted unit: the article's theorem thm:converse (source0.tex lines 491--508). The article's complete original proof of it is delivered with it (source0.tex lines 510--557, one \texttt{proof} environment of 48 lines). No mathematics is written, repaired or re-derived: the statement and the proof are the article's own lines, in the article's own notation, and no retained line is substituted or re-flowed.  No further source-local context is needed: every cross-reference of every retained line resolves inside the two delivered spans.  Nothing was omitted from any retained span, so the package carries no omission record.  No retained line contains a citation, so the wrapper carries no bibliography and no bibliography entry.  No retained line contains a figure environment, an image-inclusion command or drawing code, so no image is delivered and the package declares no asset dependency.  Editorial formatting only, in addition: an AMS \texttt{amsart} preamble that restates the article's own declarations and macros used by the retained lines, namely the environments \texttt{theorem} on separate counters, together with the standard packages those lines need.  The retained environments print in this document as Theorem 1 in order of appearance, whereas the article prints the same items with the numbering of its own full document.  The complete native source of the article as distributed in the arXiv e-print for this exact version is bundled unchanged as \texttt{sources/209442/source0.tex}.
\begin{theorem}[universal linear-reindexing criterion]\label{thm:converse}
Let $p$ be prime, $k\ge1$ and $L\in GL_k(\FF_p)$. Let $W$ be the set of
$p\times p$ integer arrays all of whose row sums and all of whose column sums
are equal to a common value $\gamma(a)$. The following are equivalent.
\begin{enumerate}
\item[\textup{(i)}] $\supp(L^{\tr}s)\ne\supp(s)$ for every
  $s\in\FF_p^k\setminus\{0\}$.
\item[\textup{(ii)}] For \emph{every} $a_1,\dots,a_k\in W$,
\[
\sum_{v\in\FF_p^k}\prod_{i=1}^{k}a_i\big((Lv)_i,\,v_i\big)
\;=\;
\sum_{v\in\FF_p^k}\prod_{i=1}^{k}a_i\big((Lv)_i,\,p-1-v_i\big)
\;=\;\prod_{i=1}^{k}\gamma(a_i).
\]
\end{enumerate}
Each of the two identities in \textup{(ii)}, taken alone, is already equivalent
to \textup{(i)}.
\end{theorem}

\begin{proof}
Let $F$ and $F'$ be the two differences (left side minus $\prod_i\gamma(a_i)$)
in (ii), regarded as functions of $(a_1,\dots,a_k)$. Both are multilinear: the
sums are multilinear by inspection, and $\prod_i\gamma(a_i)$ is a product of
linear functionals in distinct arguments.

$W$ is a free $\ZZ$-module of rank $d=(p-1)^2+1$, with coordinates the entries
$a_{x,y}$ for $x,y<p-1$ together with $\gamma(a)$: the remaining entries are
recovered by $a_{x,p-1}=\gamma-\sum_{y<p-1}a_{x,y}$,
$a_{p-1,y}=\gamma-\sum_{x<p-1}a_{x,y}$ and
$a_{p-1,p-1}=\gamma-\sum_{y<p-1}a_{p-1,y}$, and one checks that the last column
then sums to $\gamma$ automatically. The complex arrays satisfying the same
conditions form a space $W_\CC$ of dimension $d$ over $\CC$, cut out by the same
integer linear equations, so a $\ZZ$-basis of $W$ is a $\CC$-basis of $W_\CC$.
A multilinear form is determined by its values on tuples of basis elements, and
those values are the same numbers whether $F$ is read on $W^k$ or on $W_\CC^k$.
Hence $F$ vanishes on $W^k$ if and only if its $\CC$-multilinear extension
vanishes on $W_\CC^k$; in particular we may test on any $\CC$-basis of $W_\CC$,
and no question of rationality, or of the reality of an individual array, arises.

Let $\omega=e^{2\pi i/p}$, $\chi_{(s,t)}(x,y)=\omega^{sx+ty}$, and
$B=\{(0,0)\}\cup\{(s,t):s\ne0\ne t\}$. For $t\ne0$ every row sum of
$\chi_{(s,t)}$ vanishes and for $s\ne0$ every column sum vanishes, so
$\chi_{(s,t)}\in W_\CC$ for $(s,t)\in B$, while $\chi_{(0,0)}$ is the all-ones
array. Distinct characters of $(\ZZ/p)^2$ are linearly independent and
$|B|=(p-1)^2+1=d$, so these form a $\CC$-basis of $W_\CC$. Moreover
$\gamma(\chi_{(0,0)})=p$ and $\gamma(\chi_{(s,t)})=0$ for $s,t\ne0$.

Evaluate at $a_i=\chi_{(\sigma_i,\tau_i)}$ with $(\sigma_i,\tau_i)\in B$. Since
$\sum_{v\in\FF_p^k}\omega^{\langle u,v\rangle}$ is $p^k$ for $u=0$ and $0$
otherwise,
\[
\sum_{v}\prod_i\omega^{\sigma_i(Lv)_i+\tau_iv_i}
=\sum_{v}\omega^{\langle L^{\tr}\sigma+\tau,\;v\rangle}
=p^k\,\big[\,L^{\tr}\sigma+\tau=0\,\big],
\]
whereas $\prod_i\gamma(\chi_{(\sigma_i,\tau_i)})$ is $p^k$ when every
$(\sigma_i,\tau_i)=(0,0)$ and $0$ otherwise. In that first case
$\sigma=\tau=0$, both sides equal $p^k$ and $F$ vanishes. Otherwise the second
quantity is $0$, so $F$ is nonzero at this basis tuple precisely when
$\tau=-L^{\tr}\sigma$. Such a tuple exists with $(\sigma,\tau)\ne0$ exactly when
some $\sigma\ne0$ has, for every $i$, $\sigma_i$ and $(L^{\tr}\sigma)_i$ both
zero or both nonzero --- that is, $\supp(L^{\tr}\sigma)=\supp(\sigma)$. So
$F\equiv0$ if and only if (i) holds. For $F'$ the same computation acquires the
factor $\omega^{(p-1)\sum_i\tau_i}$, a root of unity and so nonzero, and the
condition $\tau=L^{\tr}\sigma$; the support condition is unchanged. Hence
$F'\equiv0$ if and only if (i) holds as well.
\end{proof}

\end{document}
